% from ega0-1-fr.tex
\oldpage[0\textsubscript{I}]{47}

que $u_x=\varphi$ (resp. $v_x=\psi$). Comme $(u\circ v)_x$ et
$(v\circ u)_x$ sont les automorphismes identiques, il existe un voisinage
ouvert $U\subset V$ de $x$ tel que $(u\circ v)|U$ et $(v\circ u)|U$ soient les
automorphismes identiques, d'où la proposition.
\end{env}

\subsection{Faisceaux cohérents}
\label{subsection:0.5.3-fr}

\begin{env}[5.3.1]
\label{0.5.3.1-fr}
On dit qu'un $\mathcal{O}_X$-Module $\mathcal{F}$ est \emph{cohérent} s'il
vérifie les deux conditions suivantes~:
\begin{enumerate}
  \item[$a)$] $\mathcal{F}$ est de type fini.
  \item[$b)$] Pour tout ouvert $U\subset X$, tout entier $n>0$, et tout
    homomorphisme
    $u:\mathcal{O}_X^n|U\to\mathcal{F}|U$, le noyau de $u$ est de type fini.
\end{enumerate}
On notera que ces deux conditions sont de caractère \emph{local}.

Pour la plupart des démonstrations des propriétés des faisceaux cohérents
rappelées dans ce qui suit, cf. (FAC), I, 2.
\end{env}

\begin{env}[5.3.2]
\label{0.5.3.2-fr}
Tout $\mathcal{O}_X$-Module cohérent admet une présentation finie (5.2.5)~; la
réciproque n'est pas nécessairement exacte, car $\mathcal{O}_X$ lui-même n'est
pas nécessairement un $\mathcal{O}_X$-Module cohérent.

Tout sous-$\mathcal{O}_X$-Module \emph{de type fini} d'un
$\mathcal{O}_X$-Module cohérent est cohérent~; toute somme directe \emph{finie}
de $\mathcal{O}_X$-Modules cohérents est un $\mathcal{O}_X$-Module cohérent.
\end{env}

\begin{env}[5.3.3]
\label{0.5.3.3-fr}
Si $0\to\mathcal{F}\to\mathcal{G}\to\mathcal{H}\to0$ est une suite exacte de
$\mathcal{O}_X$-Modules et si deux de ces $\mathcal{O}_X$-Modules sont
cohérents, il en est de même du troisième.
\end{env}

\begin{env}[5.3.4]
\label{0.5.3.4-fr}
Si $\mathcal{F}$ et $\mathcal{G}$ sont deux $\mathcal{O}_X$-Modules cohérents,
$u:\mathcal{F}\to\mathcal{G}$ un homomorphisme, $\operatorname{Im}(u)$,
$\operatorname{Ker}(u)$ et $\operatorname{Coker}(u)$ sont des
$\mathcal{O}_X$-Modules cohérents. En particulier, si $\mathcal{F}$ et
$\mathcal{G}$ sont des sous-$\mathcal{O}_X$-Modules cohérents d'un
$\mathcal{O}_X$-Module cohérent, $\mathcal{F}+\mathcal{G}$ et
$\mathcal{F}\cap\mathcal{G}$ sont cohérents.
\end{env}

\begin{env}[5.3.5]
\label{0.5.3.5-fr}
Si $\mathcal{F}$ et $\mathcal{G}$ sont deux $\mathcal{O}_X$-Modules cohérents,
il en est de même de $\mathcal{F}\otimes_{\mathcal{O}_X}\mathcal{G}$ et de
$\mathcal{H}\!om_{\mathcal{O}_X}(\mathcal{F},\mathcal{G})$.
\end{env}

\begin{env}[5.3.6]
\label{0.5.3.6-fr}
Soient $\mathcal{F}$ un $\mathcal{O}_X$-Module cohérent, $\mathcal{J}$ un
faisceau cohérent d'idéaux de $\mathcal{O}_X$. Alors le
$\mathcal{O}_X$-Module $\mathcal{J}\mathcal{F}$ est cohérent, en tant qu'image
de $\mathcal{J}\otimes_{\mathcal{O}_X}\mathcal{F}$ par l'homomorphisme
canonique
$\mathcal{J}\otimes_{\mathcal{O}_X}\mathcal{F}\to\mathcal{F}$
(5.3.4 et 5.3.5).
\end{env}

\begin{env}[5.3.7]
\label{0.5.3.7-fr}
On dit qu'une $\mathcal{O}_X$-Algèbre $\mathcal{A}$ est \emph{cohérente} si
$\mathcal{A}$ est un $\mathcal{O}_X$-Module cohérent. En particulier,
$\mathcal{O}_X$ est un \emph{faisceau cohérent d'anneaux} si, et seulement si,
pour tout ouvert $U\subset X$ et tout homomorphisme de la forme
$u:\mathcal{O}_X^p|U\to\mathcal{O}_X|U$, le noyau de $u$ est un
$(\mathcal{O}_X|U)$-Module de type fini.

Si $\mathcal{O}_X$ est un faisceau cohérent d'anneaux, tout
$\mathcal{O}_X$-Module $\mathcal{F}$ admettant une présentation finie (5.2.5)
est cohérent, en vertu de (5.3.4).

L'\emph{annulateur} d'un $\mathcal{O}_X$-Module $\mathcal{F}$ est le noyau
$\mathcal{J}$ de l'homomorphisme canonique
$\mathcal{O}_X\to\mathcal{H}\!om_{\mathcal{O}_X}(\mathcal{F},\mathcal{F})$
qui, à toute section $s\in\Gamma(U,\mathcal{O}_X)$ fait correspondre la
multiplication par $s$ dans
$\operatorname{Hom}(\mathcal{F}|U,\mathcal{F}|U)$~; si $\mathcal{O}_X$ est
cohérent et si $\mathcal{F}$ est un $\mathcal{O}_X$-Module cohérent,
$\mathcal{J}$ est cohérent (5.3.4 et 5.3.5) et pour tout $x\in X$,
$\mathcal{J}_x$ est l'annulateur de $\mathcal{F}_x$ (5.2.6).
\end{env}

